\chapter{Бимодульная совместимость повышения канального триплета}
\label{ch:v8-baryon-c0-channel-triplet-bimodule-compatibility}

\section{Физическое разложение каналов}

Условный Hodge-интертвейнер предыдущего гейта требовал считать три столбца
coherence-матрицы одним стандартным триплетом. Однако в конечной геометрии
они имеют типы
\begin{equation}
 W_e,W_X=(\mathbb C,\mathbb C,R),
 \qquad
 W_Y=(\mathbb H,\mathbb C,R),
 \qquad
 W=W_{eX}\oplus W_Y\simeq\mathbb C^2\oplus\mathbb C.
 \label{eq:v8-channel-promotion-current-types}
\end{equation}
Тем самым центральный проектор типа можно выбрать в виде
\begin{equation}
 P_Y=\operatorname{diag}(0,0,1).
 \label{eq:v8-channel-promotion-type-projector}
\end{equation}
Физическая замена базиса обязана коммутировать с $P_Y$; сохранение одной
нормы на случайном пространстве коэффициентов этого условия не заменяет
\cite{v8-channel-promotion-krajewski,v8-channel-promotion-jureit}.

\section{Коммутант типа и непрерывный запрет}

Для общей матрицы $A\in M_3(\mathbb C)$ точная система даёт
\begin{equation}
 [A,P_Y]=0
 \quad\Longleftrightarrow\quad
 A=\begin{pmatrix}A_{2\times2}&0\\0&a_Y\end{pmatrix},
 \qquad
 \rank\mathcal C_P=4,
 \qquad
 \dim\ker\mathcal C_P=5.
 \label{eq:v8-channel-promotion-complex-commutant}
\end{equation}
Следовательно, бимодульный коммутант равен
\begin{equation}
 \operatorname{End}_{\mathcal A-\mathcal A}(W)
 \simeq M_2(\mathbb C)\oplus\mathbb C,
 \label{eq:v8-channel-promotion-commutant-algebra}
\end{equation}
а вещественный ортогональный коммутант имеет лишь
\begin{equation}
 \mathfrak{so}(3)\cap\{P_Y\}'
 =\mathbb R L_{12}\simeq\mathfrak{so}(2),
 \qquad
 \dim=1.
 \label{eq:v8-channel-promotion-orthogonal-commutant}
\end{equation}
Два генератора, необходимые для смешивания $Y$ с $e,X$, запрещены:
\begin{equation}
 [L_{12},P_Y]=0,
 \qquad
 [L_{13},P_Y]\ne0,
 \qquad
 [L_{23},P_Y]\ne0.
 \label{eq:v8-channel-promotion-so3-generator-test}
\end{equation}
Поэтому честное стандартное действие $SO(3)$ на текущем $W$ не существует.

\section{Дискретная и spin-cover лазейки}

Стандартный тетраэдрический триплет также не сохраняет разложение $2+1$.
Для циклической перестановки $R_3$ и полуоборота $R_2$ имеем
\begin{equation}
 [R_3,P_Y]\ne0,
 \qquad
 [R_2,P_Y]=0.
 \label{eq:v8-channel-promotion-a4-type-test}
\end{equation}
Значит, полный $A_4$ нельзя использовать как замену непрерывного повышения.

На инфинитезимальном уровне можно вложить $\mathfrak{so}(3)\simeq
\mathfrak{su}(2)$ в $\mathfrak u(2)$, но фундаментальное комплексное
представление является представлением $Spin(3)=SU(2)$, а не $SO(3)$.
Центр накрытия действует как
\begin{equation}
 z_{2\pi}=\operatorname{diag}(-1,-1,1)\ne I_3,
 \qquad
 z_{2\pi}^2=I_3.
 \label{eq:v8-channel-promotion-spin-cover-obstruction}
\end{equation}
Следовательно, оно не спускается к требуемому групповому действию и в любом
случае оставляет приводимый тип $\mathbf2\oplus\mathbf1$, а не векторный
$\mathbf3$.

\section{Минимальная неразрушающая условная архитектура}

Не меняя старый $Y_R$, можно добавить один новый endpoint того же типа, что
$e_R,X_R$:
\begin{equation}
 Z_R=(\mathbb C,\mathbb C,R),
 \qquad
 W_{\rm ext}=W_{eXZ}\oplus W_Y\simeq\mathbb C^3\oplus\mathbb C.
 \label{eq:v8-channel-promotion-minimal-extension}
\end{equation}
Тогда проектор типа становится
\begin{equation}
 P_Y^{\rm ext}=\operatorname{diag}(0,0,0,1),
 \qquad
 \{P_Y^{\rm ext}\}'\simeq M_3(\mathbb C)\oplus\mathbb C,
 \qquad
 \dim_{\mathbb C}=10.
 \label{eq:v8-channel-promotion-extended-commutant}
\end{equation}
Изотипический проектор $P_{eXZ}=I_4-P_Y^{\rm ext}$ канонически выделяет
трёхмерный блок, на котором стандартный $SO(3)$ уже совместим с
бимодульными типами. Для двух строк coherence-матрицы выбранная подцепь
сохраняет размерности
\begin{equation}
 \dim H_0=1,
 \qquad
 \dim H_1=2\cdot3=6,
 \qquad
 \dim H_2=\binom22\binom32=3.
 \label{eq:v8-channel-promotion-selected-chain-dimensions}
\end{equation}
Но полный расширенный стрелочный носитель имеет размерность
\begin{equation}
 \dim\operatorname{Hom}(\mathbb C^4,\mathbb C^2)=8,
 \qquad
 \Delta\dim_{\mathbb C}=2.
 \label{eq:v8-channel-promotion-extension-cost}
\end{equation}
Старый $Y$ становится spectator-каналом, а прежний rank-one конденсат на
$(e,X,Y)$ не наследуется автоматически подпространством $(e,X,Z)$. Нужны
новые endpoint, две стрелки, новый конденсат и диагональный family--channel
lock.

Реестры принимают вид
\begin{equation}
 \begin{aligned}
 N_{\rm current\ promotion}&=0/5,
 &N_{\rm extension\ shape}&=5/5,\\
 N_{\rm new\ origin}&=0/4,
 &\mathcal M_{\rm new}&=\{Z_R,\ 2\ arrows,\ B_{eXZ,*},\ SO(3)_{\rm diag}\}.
 \end{aligned}
 \label{eq:v8-channel-promotion-ledgers}
\end{equation}

\begin{theorem}[Бимодульный запрет текущего повышения]
На текущем канальном носителе всякая допустимая замена базиса сохраняет
$W_{eX}\oplus W_Y$. Поэтому ни стандартный $SO(3)$, ни стандартный $A_4$
не могут смешивать три coherence-канала в один физический триплет.
Spin-cover на двухмерном блоке не устраняет препятствие.
\end{theorem}

\begin{proposition}[Минимальная условная изотипическая лазейка]
Добавление одного endpoint $Z_R=(\mathbb C,\mathbb C,R)$ создаёт канонически
выделенный изотипический блок $W_{eXZ}\simeq\mathbb C^3$, допускающий
стандартный $SO(3)$ и подцепь размерностей $1$--$6$--$3$. Эта архитектура
не наследует старый конденсат и содержит четыре новых родительских входа.
\end{proposition}

\begin{nogocriterion}[Запрет смешивания неэквивалентных endpoint]
Нельзя повышать случайную $U(3)$-симметрию нормы стрелочных коэффициентов до
физической симметрии, если она не коммутирует с центральными проекторами
бимодульных типов. Ни совпадение размерностей, ни переход к $A_4$, ни
инфинитезимальное $\mathfrak{su}(2)$ не отменяют это требование.
\end{nogocriterion}

\begin{protocol}\begin{enumerate}
\item текущие типы каналов: \textbf{$2(\mathbb C,\mathbb C,R)\oplus(\mathbb H,\mathbb C,R)$};
\item бимодульный коммутант: \textbf{$M_2(\mathbb C)\oplus\mathbb C$};
\item допустимая ортогональная алгебра: \textbf{$\mathfrak{so}(2)$};
\item стандартный $SO(3)$ совместим: \textbf{нет};
\item стандартный $A_4$ совместим: \textbf{нет};
\item spin-cover исправляет запрет: \textbf{нет};
\item текущий promotion-ledger: \textbf{$0/5$};
\item минимальное расширение: \textbf{один новый $Z_R=(\mathbb C,\mathbb C,R)$};
\item расширенный коммутант: \textbf{$M_3(\mathbb C)\oplus\mathbb C$};
\item выбранная coherence-подцепь: \textbf{$1$--$6$--$3$};
\item условная форма расширения: \textbf{$5/5$};
\item новых parent-входов: \textbf{четыре};
\item portal parent-origin: \textbf{не закрыт};
\item следующий гейт: \textbf{минимальное изотипическое канальное расширение}.
\end{enumerate}\end{protocol}

Машинный сертификат сохранён в
\path{s2t_v8_baryon_c0_singlet_triplet_central_gap_coherence_channel_triplet_promotion_bimodule_compatibility_gate_results.json}.

\begin{thebibliography}{9}
\bibitem{v8-channel-promotion-krajewski}
T.~Krajewski,
\emph{Classification of Finite Spectral Triples},
arXiv:hep-th/9701081.
\bibitem{v8-channel-promotion-jureit}
J.-H.~Jureit, C.~A.~Stephan,
\emph{On a Classification of Irreducible Almost-Commutative Geometries IV},
arXiv:hep-th/0610040.
\end{thebibliography}